% Cote 119: offprints, reports, courses, CV and letters (typed and printed part): archive cote 119, site page 292 (PDF page 293).
% typescript. Draft compared: none.
% AI-transcribed by Claude Opus 5.5 (claude-opus-5-5, adaptive thinking, effort high), Message Batches API, batch r119_full (msgbatch_01DwZNm5Y7YsKBWyAN4sFs8w), 2026-11-13; not yet checked by eye.
\scanpage{119/292}
\mspage{- 5 -}

3.3. Le groupe $G_2$

Supposons que $G$ soit égal au groupe exceptionnel $G_2$ , et soit déployé.
On sait qu'il y a alors des bijections naturelles entre les trois ensembles
suivants :

$H^1(k, G_2)$ ;

classes d'algèbres d'octonions sur $k$;

classes de 3-formes de Pfister sur $k$.

Il résulte de là, et des théorèmes cités ci-dessus, que, si l'on prend
pour $\hand{$\rho$}$ la représentation fondamentale de degré $7$ de $G_2$ , l'application
$i_{\hand{$\rho$}}$ correspondante est une \uline{bijection de} $H^1(k, G_2)$ \uline{sur le sous-ensemble}
\uline{de} $H^3(k, Z/2Z)$ \uline{formé des éléments décomposables} (i.e. cup-produits
de trois éléments de $H^1(k, Z/2Z)$). Cela donne une description cohomologique
tout à fait satisfaisante de l'ensemble $H^1(k, G_2)$.

On peut aller un peu plus loin. Notons $i$ l'injection de $H^1(k, G_2)$
dans $H^3(k, Z/2Z)$ que nous venons de définir. Soit $\hand{$\rho$}$ une représentation
irréductible quelconque de $G_2$ ; il lui correspond d'après (3.2.1) une
application
\[ i_{\hand{$\rho$}} : H^1(k, G_2) \quad -> \quad H^3(k, Z/2Z). \]
On désire comparer $i_{\hand{$\rho$}}$ à $i$. Le résultat est le suivant (je me borne
ici au cas où le corps de base est de caractéristique $0$) :

(3.3.1) \uline{On a}, \uline{soit} $i_{\hand{$\rho$}} = i$, \uline{soit} $i_{\hand{$\rho$}} = 0$ .

De façon plus précise, soit $m_1\hand{$\omega$}_1 + m_2\hand{$\omega$}_2$ le poids dominant de $\hand{$\rho$}$ , écrit
comme combinaison linéaire des poids fondamentaux $\hand{$\omega$}_1$ et $\hand{$\omega$}_2$ ($\hand{$\omega$}_1$ corres-%
pondant à la représentation de degré $7$, et $\hand{$\omega$}_2$ à la représentation adjointe).
On peut déterminer (grâce à des formules qui m'ont été communiquées par
J.Tits) dans quel cas on a $i_{\hand{$\rho$}} = i$ ; on trouve que cela se produit si
et seulement si le couple $(m_1, m_2)$ est congru (mod 8) à l'un des douze
couples suivants :

(0,2), (0,3), (1,0), (1,4), (2,0), (2,3), (4,3), (4,6), (5,2), (5,6),
(6,3), (6,4).

Ainsi, pour la représentation adjointe\hand{,}qui correspond à (0,1), on
a $i_{\hand{$\rho$}} = 0$. On peut préciser ceci en déterminant explicitement la forme
de Killing $\mathrm{Kill}_x$ de la k-forme de $G_2$ associée à un élément donné
$x \hand{$\in$} H^1(k, G_2)$. Si $q_x \hand{$=$} \hand{$\langle 1 \rangle \oplus$} q_x^{\hand{$\circ$}}$ est la 3-forme de Pfister associée à $x$
(i.e. la \uline{forme norme} de l'algèbre d'octonions correspondante), on trouve
que $\mathrm{Kill}_x$ \uline{est isomorphe à} $\hand{$\langle -1,-3 \rangle \otimes$} q_x^{\hand{$\circ$}}$ .

\margin{292}
